\subsection{中微子的物质传播与 MSW 共振}
\label{sec:sm-neutrino-matter-applications}

真空振荡比较不同质量分量的传播相位；弱相互作用还允许介质在不记录中微子路径的相干前向散射中改变这些相位。因而一次可分辨碰撞的概率很小，并不意味着长距离传播时介质的作用可以忽略。以下在已引入中微子质量的标准模型有效理论中，把前面的低能弱相互作用与味空间演化联系起来。通常的三活性味物质效应与太阳中微子应用可参阅综述\cite{GonzalezGarciaYokoyama2024Neutrino}。

设 $z$ 为沿束流的距离，$\boldsymbol\psi_f(z)=(\psi_e,\psi_\mu,\psi_\tau)^{\mathsf T}$ 为归一化味振幅列向量。对固定主导能量 $E_\nu$，先从质量壳关系展开 $p_ic=E_\nu-m_i^2c^4/(2E_\nu)+\mathcal O(m_i^4c^8/E_\nu^3)$，再在局域前向平均场中传播，得到
\begin{equation}
 \begin{aligned}
 \ii\hbar c\frac{\dd\boldsymbol\psi_f}{\dd z}
 &=H_f(z)\boldsymbol\psi_f(z),\\
 H_f(z)&=U\frac{\operatorname{diag}(m_1^2,m_2^2,m_3^2)c^4}{2E_\nu}U^\dagger
       +\operatorname{diag}(V_{\rm CC}(z),0,0),\\
 V_{\rm CC}(z)&=\sqrt2G_F^{(E)}(\hbar c)^3 n_e(z).
 \end{aligned}
 \label{eq:sm-matter-flavor-hamiltonian}
\end{equation}
这里 $U=U_{\rm PMNS}$；$n_e$ 是无极化、静止普通物质的电子数密度，单位为 $\mathrm{m^{-3}}$；$G_F^{(E)}$ 的量纲是能量的负二次方，故 $V_{\rm CC}$ 是能量。中性流在三个活性味中给出相同的树级势，已作为整体相位扣除。有惰性味、极化介质或中微子自身的显著前向势时，必须重新构造相应矩阵。反中微子在同一普通物质中满足 $U\to U^*$、$V_{\rm CC}\to -V_{\rm CC}$。

这一步保留超相对论色散的首个质量修正与弱相互作用前向振幅的最低阶。它要求 $m_i^2c^4/E_\nu^2\ll1$、局域弱作用的动量不变量远低于 $m_W^2c^2$、介质变化长度远大于 $\hbar c/E_\nu$，并忽略传播段内可分辨碰撞造成的吸收和非相干散射。质量波包还需保持重叠。密度变化慢于德布罗意波长只保证局域传播方程成立，并不自动保证味演化绝热。

对任意密度剖面，这个有效方程的形式解是
\begin{equation}
 \boldsymbol\psi_f(L)=S_f(L,0)\boldsymbol\psi_f(0),\qquad
 S_f(L,0)=\mathcal P\exp\!\left[-\frac{\ii}{\hbar c}
 \int_0^L H_f(z)\,\dd z\right].
 \label{eq:sm-matter-path-ordered}
\end{equation}
$\mathcal P$ 把较大 $z$ 的矩阵置于左边；$P_{\alpha\to\beta}=|(S_f)_{\beta\alpha}|^2$。厄米性保证 $S_f^\dagger S_f=I$。把密度划分为常值层后，传播矩阵是各层指数按实际经过次序相乘；一般不能把势的平均值直接代入一个指数，因为不同位置的 $H_f$ 未必对易。

现在才取孤立的两味子空间 $(\nu_e,\nu_x)$，其中 $\nu_x$ 是与电子味混合的另一活性味组合。忽略第三态要求它的耦合或相关质量尺度修正小到满足所需精度；下面的二味算例不作为三味参数拟合。取实旋转
\begin{equation*}
 U_2=\begin{pmatrix}\cos\theta&\sin\theta\\-\sin\theta&\cos\theta\end{pmatrix},
 \qquad \delta_\nu=\frac{\Delta m^2c^4}{2E_\nu},\qquad
 \Delta m^2=m_2^2-m_1^2,
\end{equation*}
其中 $\delta_\nu$ 是能量差，以下选 $\Delta m^2>0$、$0<\theta<\pi/4$。扣掉迹的一半后，哈密顿量及其本征能量间隔为
\begin{equation}
 \begin{aligned}
 H_2&=\frac12\begin{pmatrix}
 V_{\rm CC}-\delta_\nu\cos2\theta&\delta_\nu\sin2\theta\\
 \delta_\nu\sin2\theta&\delta_\nu\cos2\theta-V_{\rm CC}
 \end{pmatrix},\\
 D_\nu&=\sqrt{(\delta_\nu\cos2\theta-V_{\rm CC})^2
                       +\delta_\nu^2\sin^22\theta},\\
 \cos2\theta_m&=\frac{\delta_\nu\cos2\theta-V_{\rm CC}}{D_\nu},\qquad
 \sin2\theta_m=\frac{\delta_\nu\sin2\theta}{D_\nu}.
 \end{aligned}
 \label{eq:sm-matter-two-flavor}
\end{equation}
$D_\nu$ 仍是能量，$\theta_m\in(0,\pi/2)$ 是物质中的瞬时混合角；同时给出正弦和余弦才能固定它在共振两侧的分支。常密度下的转换概率为
\begin{equation}
 P_{e\to x}(L)=\sin^22\theta_m\,
 \sin^2\!\left(\frac{D_\nu L}{2\hbar c}\right),\qquad
 P_{e\to e}(L)=1-P_{e\to x}(L).
 \label{eq:sm-matter-constant-probability}
\end{equation}
共振条件 $V_{\rm CC}=\delta_\nu\cos2\theta$ 使混合振幅达到一；能级仍隔着 $D_\nu=\delta_\nu\sin2\theta$，所以这是避免交叉。在所选质量差与角度分支下，普通物质中的反中微子没有这一共振。$V_{\rm CC}\to0$ 恢复真空相位 $\Delta m^2c^3L/(4E_\nu\hbar)$；$|V_{\rm CC}|\gg\delta_\nu$ 时常密度振幅被抑制。共振处再令 $\theta\to0$，所需振荡长度发散：任何固定有限 $L$ 上的转换仍趋于零，不能只看最大混合振幅就断言完全转换。

\begin{exbox}[由数密度计算共振能量和传播长度]
取教学输入 $\Delta m^2c^4=2.50\times10^{-3}\,\mathrm{eV^2}$、$\theta=8.6^\circ$、$n_e=1.50\times10^{30}\,\mathrm{m^{-3}}$。使用
\begin{equation*}
 G_F^{(E)}=1.16638\times10^{-23}\,\mathrm{eV^{-2}},\qquad
 \hbar c=1.973269804\times10^{-7}\,\mathrm{eV\,m},
\end{equation*}
有 $V_{\rm CC}=1.90\times10^{-13}\,\mathrm{eV}$。直接代入共振条件，
\begin{equation*}
 E_{\rm res}=\frac{\Delta m^2c^4\cos2\theta}{2V_{\rm CC}}
 \simeq\SI{6.28}{\giga\electronvolt},\qquad
 L_{\rm max}=\frac{\pi\hbar c}{D_\nu(E_{\rm res})}
 \simeq1.05\times10^4\,\mathrm{km}.
\end{equation*}
$L_{\rm max}$ 是第一次完全转换的距离，不是完整周期；完整周期为 $2L_{\rm max}$。同样的密度和能量下，反中微子的混合振幅仅约 $0.0234$。这些数值说明长基线中的物质效应为何可观，但实际穿过地球的计算应使用分层密度和三味演化。
\end{exbox}

\figref{fig:sm-neutrino-matter} 将能量依赖的混合振幅与给定能量下的距离依赖分开显示。

\begin{figure}[htbp]
 \centering
 \includegraphics[width=0.92\linewidth]{sm_neutrino_matter.pdf}
 \caption{二味常密度传播的教学计算。左图用同一普通物质的电子数密度比较中微子和反中微子的有效混合振幅；右图固定在中微子的共振能量，比较物质与真空的转换概率。大混合角还需要足够传播距离才能变成大转换概率。参数见正文算例，图中没有加入第三味、密度梯度或实验能量分辨率。}
 \label{fig:sm-neutrino-matter}
\end{figure}

对随位置变化的密度，令 $R_m(z)$ 为与 $U_2$ 同形式、角度为 $\theta_m(z)$ 的实旋转，定义瞬时本征基振幅 $\boldsymbol\psi_m=R_m^\dagger\boldsymbol\psi_f$。运动方程变为
\begin{equation}
 \begin{aligned}
 \ii\hbar c\frac{\dd\boldsymbol\psi_m}{\dd z}
 &=\left[\operatorname{diag}(-D_\nu/2,D_\nu/2)
        -\ii\hbar c R_m^\dagger\frac{\dd R_m}{\dd z}\right]\boldsymbol\psi_m,\\
 \epsilon_{\rm ad}(z)&=\frac{\hbar c|\dd\theta_m/\dd z|}{D_\nu}
 =\frac{\hbar c|\delta_\nu\sin2\theta\,\dd V_{\rm CC}/\dd z|}{2D_\nu^3}.
 \end{aligned}
 \label{eq:sm-matter-adiabatic-condition}
\end{equation}
无能隙闭合、平滑单次密度变化时，$\epsilon_{\rm ad}\ll1$ 是抑制瞬时本征态间跃迁的局域判据。突变界面或很长的周期调制应直接检查传播矩阵，不能仅用某一点的比值判断累计跃迁。在共振处，判据变成 $\hbar c|V_{\rm CC}'|/(2\delta_\nu^2\sin^22\theta)\ll1$。密度越平缓，态越能沿避免交叉的同一条本征分支传播。

若电子味在 $z=0$ 产生，随后绝热传播到真空，且探测的能量或路径平均已洗掉两条分支的相对相位，存活概率为
\begin{equation}
 \overline P_{e\to e}
 =\cos^2\theta_m(0)\cos^2\theta+\sin^2\theta_m(0)\sin^2\theta
 =\frac12[1+\cos2\theta_m(0)\cos2\theta].
 \label{eq:sm-matter-averaged-survival}
\end{equation}
低密度产生时得到真空平均 $1-\tfrac12\sin^22\theta$；高密度产生时 $\theta_m(0)\to\pi/2$，得到 $\sin^2\theta$。这里改变的是本征态随密度的组成，因此高密度出发的绝热转换与一直待在高密度介质中的小振幅振荡并不矛盾。对相干、单能且已知基线的束流，应保留相位干涉项，而不能使用这一平均公式。

\begin{exbox}[指数剖面的绝热性检查]
取 $\Delta m^2c^4=7.50\times10^{-5}\,\mathrm{eV^2}$、$\theta=33.4^\circ$、$E_\nu=\SI{10}{\mega\electronvolt}$，并设理想化密度 $n_e(z)=n_{e0}\ee^{-z/r_n}$，其中 $n_{e0}=6.00\times10^{31}\,\mathrm{m^{-3}}$、$r_n=6.60\times10^7\,\mathrm m$。这只是用于练习的平滑介质，并未包含真实产生区分布。
在起点有 $V_{\rm CC}(0)\simeq7.60\times10^{-12}\,\mathrm{eV}$，$\delta_\nu=3.75\times10^{-12}\,\mathrm{eV}$。共振位置由
\begin{equation*}
 z_{\rm res}=r_n\ln\frac{V_{\rm CC}(0)}{\delta_\nu\cos2\theta}
 \simeq1.08\times10^8\,\mathrm m
\end{equation*}
确定。利用 $|V_{\rm CC}'|_{\rm res}=\delta_\nu\cos2\theta/r_n$ 得
\begin{equation*}
 \epsilon_{\rm ad}(z_{\rm res})
 =\frac{\hbar c\cos2\theta}{2r_n\delta_\nu\sin^22\theta}
 \simeq1.86\times10^{-4}.
\end{equation*}
这一平滑单次剖面允许绝热近似；相位平均后 $\overline P_{e\to e}\simeq0.328$。它尚未达到无限高起始密度的极限 $\sin^2\theta\simeq0.303$，而真空平均约为 $0.578$。使用具体初始混合角，可以避免过早套用高密度极限产生的误差。
\end{exbox}

\begin{derivation}[从带电流平均场到物质势]
为跟踪量纲，取规范归一化空间场 $\psi$ 的量纲为 $\mathrm{m^{-3/2}}$。低能带电流经 Fierz 重排后，其哈密顿密度为
\begin{equation*}
 \mathcal H_{\rm CC}^{\rm eff}
 =\frac{G_F^{(E)}(\hbar c)^3}{\sqrt2}
 [\bar\nu_e\gamma^\mu(1-\gamma^5)\nu_e]
 [\bar e\gamma_\mu(1-\gamma^5)e].
\end{equation*}
系数的单位是能量乘体积，两个流的乘积单位是 $\mathrm{m^{-6}}$，所以总量为能量密度。介质静止、无极化且忽略正电子时，
\begin{equation*}
 \langle\bar e\gamma^\mu e\rangle=(n_e,\boldsymbol0),\qquad
 \langle\bar e\gamma^\mu\gamma^5e\rangle=0.
\end{equation*}
对传播的左手中微子，$1-\gamma^5=2P_L$，因而
\begin{equation*}
 \langle\mathcal H_{\rm CC}^{\rm eff}\rangle_e
 =\sqrt2G_F^{(E)}(\hbar c)^3n_e\,\nu_{eL}^\dagger\nu_{eL}.
\end{equation*}
这就确定了\eqnrefs{eq:sm-matter-flavor-hamiltonian}的势及正号。有正电子时树级矢量流数密度改为净密度 $n_{e^-}-n_{e^+}$；反中微子的相对势反号。

味基与质量基的方向也可直接检查。设 $\ell_L^{\rm weak}=U_{eL}\ell_L$、$\nu_L^{\rm weak}=U_\nu\nu_L$，则
\begin{equation*}
 \bar\nu_L^{\rm weak}\gamma^\mu\ell_L^{\rm weak}
 =\bar\nu_{iL}\gamma^\mu(U_\nu^\dagger U_{eL})_{i\alpha}\ell_{\alpha L}
 =\bar\nu_{iL}\gamma^\mu U_{\alpha i}^*\ell_{\alpha L}.
\end{equation*}
故产生味态的质量振幅是 $U_{\alpha i}^*$，探测为 $\beta$ 味的投影是 $U_{\beta i}$；组合成 $U\operatorname{diag}(m_i^2)U^\dagger$ 与真空振幅的顺序一致。
\end{derivation}

\begin{derivation}[二味指数与瞬时基中的附加耦合]
将 $U_2\operatorname{diag}(0,\delta_\nu)U_2^{\mathsf T}+\operatorname{diag}(V_{\rm CC},0)$ 减去 $(\delta_\nu+V_{\rm CC})I/2$，得到\eqnrefs{eq:sm-matter-two-flavor}。用 Pauli 矩阵写成
\begin{equation*}
 H_2=\frac12[\delta_\nu\sin2\theta\,\sigma_x+
 (V_{\rm CC}-\delta_\nu\cos2\theta)\sigma_z],\qquad
 H_2^2=\frac{D_\nu^2}{4}I.
\end{equation*}
将指数级数的偶数项与奇数项分开，
\begin{equation*}
 \ee^{-\ii H_2L/(\hbar c)}
 =I\cos\frac{D_\nu L}{2\hbar c}
 -\frac{2\ii H_2}{D_\nu}\sin\frac{D_\nu L}{2\hbar c}.
\end{equation*}
非对角元为 $-\ii(\delta_\nu\sin2\theta/D_\nu)\sin[D_\nu L/(2\hbar c)]$，取模平方就是\eqnrefs{eq:sm-matter-constant-probability}。

密度变化时把 $\boldsymbol\psi_f=R_m\boldsymbol\psi_m$ 代回运动方程。乘以 $R_m^\dagger$ 并展开乘积导数，得
\begin{equation*}
 R_m^\dagger R_m'=\theta_m'\begin{pmatrix}0&1\\-1&0\end{pmatrix},\qquad
 \theta_m'=\frac{\delta_\nu\sin2\theta}{2D_\nu^2}V_{\rm CC}'.
\end{equation*}
第二式可由 $2\theta_m=\operatorname{atan2}(\delta_\nu\sin2\theta,\delta_\nu\cos2\theta-V_{\rm CC})$ 求导得到。基底导数的非对角耦合大小为 $\hbar c|\theta_m'|$，与能隙比较后得到\eqnrefs{eq:sm-matter-adiabatic-condition}。绝热时两条分支只积累相位，初始权重分别为 $\cos^2\theta_m(0)$ 和 $\sin^2\theta_m(0)$；投影到真空电子味再丢弃已平均为零的交叉项，就得到\eqnrefs{eq:sm-matter-averaged-survival}。
\end{derivation}
